\documentclass[a4paper, 12pt]{article}

% Добавляем новые пакеты
\usepackage[english, russian]{babel}    % Пакет для поддержки русского языка
\usepackage[T1, T2A]{fontenc}           % Загружаем пакет с поддержкой кирилицы
\usepackage[utf8]{inputenc}             % Загружаем пакет с поддержкой кодировки файла и символов
\usepackage{graphicx}                   % Загружаем пакет для добавления различной графики (pdf, jpg, png)
\usepackage{import}                     % Загружаем пакет для продвинутого импорта, в частности изображений pdf_tex из графического редактора inkscape (\import{<path to file>}{<filename>.pdf_tex})
\usepackage{subcaption}                 % Пакет для поддержки сетки из фигур, например в две колонки или матрицей
\usepackage{float}                      % Пакет для улучшенного размещения изображенией (ниже используется опция H в команде \begin{figure}[H])
\usepackage{geometry}                   % Пакет для управления отступами на странице (см. далее вызываем \geometry)
\usepackage{amsmath}                    % Добавляем крутую поддержку математики, в частности /align
\usepackage{amsfonts}                   % Пкет с мтематическими шрифтами вроде \mathbb и др.
\usepackage{tabu}                       % Пакет с прокаченными таблицами
\usepackage{booktabs}                   % Крутые вертикальные разделители
\usepackage[dvipsnames]{xcolor}         % Рассширенная поддержка цветов
\usepackage[colorlinks]{hyperref}
\usepackage{indentfirst}                % Пакет, чтобы починить красную строку
\usepackage{listings}                   % Поддержка листингов, оформления кода программы
\usepackage{vpdtitle}               % Титульный лист (смотри файл vpdtitle.sty)


% Устанавливает отступы на странице
\geometry{left=2cm, right=2cm, bottom=2cm, top=2cm}

% Настараиваем красивые цвета для ссылок на формулы, фигуры (изображения) и другие элементы.
\hypersetup{
    linkcolor=blue,
    filecolor=magenta,
    urlcolor=cyan
}

% Настраиваем супер красивые листинги с попомщью пакетов listings и xcolor
\definecolor{codegreen}{rgb}{0,0.6,0}
\definecolor{codegray}{rgb}{0.5,0.5,0.5}
\definecolor{codepurple}{rgb}{0.58,0,0.82}
\definecolor{backcolour}{rgb}{0.95,0.95,0.92}

\lstdefinestyle{codestyle}{
    backgroundcolor=\color{backcolour},
    commentstyle=\color{codegreen},
    keywordstyle=\color{magenta},
    numberstyle=\tiny\color{codegray},
    stringstyle=\color{codepurple},
    basicstyle=\ttfamily\footnotesize,
    breakatwhitespace=false,
    breaklines=true,
    captionpos=b,
    keepspaces=true,
    numbers=left,
    numbersep=5pt,
    showspaces=false,
    showstringspaces=false,
    showtabs=false,
    tabsize=2
}

\lstset{style=codestyle}
\lstset{extendedchars=\true}

%%%%%%%%%%%%%%%%%%%%%%%%%%%%%%%%%%%%%%%%%%%%%%%%
%% Заполнение титульного листа
%%%%%%%%%%%%%%%%%%%%%%%%%%%%%%%%%%%%%%%%%%%%%%%%
\title{
Отчёт о лабораторной работе по Git и GitHub
}

% Тут укажите автора
\author{Фам Ба Дат}
% Тут Номер группы, инфомация для связи (телеграм или почта)
\affil{407879, @BaDatPham}

% Впишите данные о том, кто будет проверять работу
\supervisor{Василий А. Николаевич}

% Краткое описание того, что сделали в работе
\abstract{
В работе рассмотрено использование Git и GitHub для управления версиями.
Описаны результаты выполнения трёх заданий: создание репозитория,
работа с командной строкой и настройка \texttt{.gitignore}.
}
% Подберите 3-5 ключевых слова, отличающих эту работу от всех остальных
\keywords{
    Git; GitHub; LaTeX; контроль версий
}

\date{\today}




\begin{document}

\maketitle

\tableofcontents
\newpage

\section{Введение}
В данном отчёте представлены результаты выполнения заданий 1, 2 и 3,
включая создание репозитория, работу с Git из командной строки и управление файлами
с помощью \texttt{.gitignore}. Отчёт подготовлен с использованием системы
вёрстки LaTeX.

\section{Задание 1: Создание файла goals.md}

\subsection{Содержимое файла goals.md}
\begin{verbatim}
# Мой опыт
Я имею около года опыта программирования на Python и C++ в рамках университетских курсов. 
Опыта командной работы с Git у меня пока немного, но я хочу это исправить в рамках данного курса.
[Мой Github](https://github.com/PhamBaDat)

# Мои цели
1. Освоить рабочий процесс с Git и GitHub в разработке ПО.
2. Улучшить навыки программирования и решение задач.
3. Успешно выполнить все задания и проекты курса.

# Мои задачи
- Выполнять все выданные задания.
- Регулярно делать commit и push кода на GitHub.
- Изучать дополнительную литературу по Git, Markdown.
- Участвовать в обсуждениях и задавать вопросы преподавателю.
- Уделять не менее 1 часа в день практике программирования.

![Иллюстрация](resources/avatar.png)
\end{verbatim}

\subsection{Результат}
Файл \texttt{goals.md} был успешно создан на GitHub и содержит разделы
«Мой опыт», «Мои цели», «Мои задачи», а также изображение в папке \texttt{resources/}.

\section{Задание 2: Работа с Git из командной строки}

\subsection{Содержимое файла info.md}
\begin{verbatim}
Pham Ba Dat
Window
20-09-2026 10:00
\end{verbatim}

\subsection{Содержимое файла .gitignore}
\begin{verbatim}
data/
\end{verbatim}

\subsection{Объяснение результата попытки коммита data/dataset.csv}
При выполнении команды \texttt{git add data/dataset.csv} Git выдал предупреждение
о том, что этот файл игнорируется файлом \texttt{.gitignore}:

\begin{verbatim}
The following paths are ignored by one of your .gitignore files:
data/dataset.csv
hint: Use -f if you really want to add them.
\end{verbatim}

Таким образом, файл \texttt{dataset.csv} не был добавлен в область подготовленных
файлов (staging area) и не мог быть закоммичен. Это доказывает, что
\texttt{.gitignore} работает корректно: он предотвращает отслеживание Git
нежелательных файлов и папок.

\section{Заключение}
В ходе выполнения заданий я освоил работу с Git и GitHub как через веб-интерфейс,
так и через командную строку, понял принцип работы \texttt{.gitignore}
и научился оформлять отчёты с помощью LaTeX.

\end{document}